\documentclass[11pt]{article}

\usepackage[T1]{fontenc}
\usepackage[utf8]{inputenc}
\usepackage{amsmath,amssymb,amsthm,mathtools}
\usepackage{booktabs}
\usepackage{array}
\usepackage[margin=1in]{geometry}
\usepackage{microtype}
\usepackage[hidelinks]{hyperref}
\hypersetup{
  pdftitle={Star-Pack Dominance and Forced-Prime CRT Lower Bounds for the Agoh--Giuga Conjecture},
  pdfauthor={Anonymous}
}

\newtheorem{theorem}{Theorem}[section]
\newtheorem{proposition}[theorem]{Proposition}
\newtheorem{lemma}[theorem]{Lemma}
\newtheorem{corollary}[theorem]{Corollary}
\theoremstyle{remark}
\newtheorem{remark}[theorem]{Remark}

\newcommand{\PP}{\mathcal{P}}
\newcommand{\Podd}{\mathbb{P}_{\mathrm{odd}}}
\newcommand{\capf}{\operatorname{cap}}
\newcommand{\gap}{\operatorname{gap}}

\title{Star-Pack Dominance and Forced-Prime CRT Lower Bounds\\
for the Agoh--Giuga Conjecture}
\author{Anonymous}
\date{}

\begin{document}
\maketitle

\begin{abstract}
The Agoh--Giuga conjecture asserts that no composite integer $n$ satisfies
$\sum_{j=1}^{n-1} j^{n-1} \equiv -1 \pmod n$; it remains open. The strongest
prior published computational lower bound, due to Borwein, Maitland, and
Skerritt (2013), used ten days of tree trimming to depth $311$ to prove that any
counterexample has at least $4{,}771$ distinct prime factors and exceeds
$10^{19{,}907}$, while noting that normal-set tree pruning can never exceed the
$8{,}135$-prime threshold of the greedy normal set $L_{27,692}$.

We first prove an elementary four-line \emph{Star-Pack Reciprocal Dominance
Lemma} and combine it with an exhaustive $7$-leaf partition of the odd primes to
establish the \emph{Exact Extremal Normal-Set Theorem} without external
large-scale search: every finite normal set $S$ of odd primes with
$\sum_{p\in S} 1/p > 1$ satisfies $|S| \ge 8{,}135$ (not uniquely attained),
$\max(S) \ge 321{,}053$ (equality iff $S = L_{27,692}$), and
$\prod_{p\in S} p \ge P(L_{27,692})$ ($40{,}733$ decimal digits, equality iff
$S = L_{27,692}$). This exhausts the Borwein--Maitland--Skerritt algorithm
($j_d = 8{,}135$ for all $d \ge 27{,}692$) and establishes that outside the
extremal leaf $\{5,7,23,83\}\subseteq S$, $3\notin S$, the exact minimum of
$\max(S)$ is $352{,}757$.

We then break past the normal-set ceiling by coupling an elementary
$(\mathbb{Z}/16\mathbb{Z})^\times$ congruence (giving $m - a \equiv 1 \pmod 8$,
independently obtained via unimodular lattices by Mavrotheris in 2026), a
three-line elementary splitting bound for Giuga values $a \ge 2$
($\omega(n)\ge 329{,}188$, $P(n)\ge 1{,}284{,}211$, $\ge 1{,}808{,}611$ digits),
and forced-prime Chinese Remainder Theorem congruences refuted by exact
$2$-dimensional lattice witnesses. Every composite counterexample to the
Agoh--Giuga conjecture therefore satisfies unconditionally
\[
  \omega(n) \ge 8{,}138,
  \qquad
  P(n) \ge 326{,}057,
  \qquad
  n \ge 10^{40{,}921} \text{ ($\ge 40{,}922$ decimal digits)},
\]
with sharper dichotomies at the smallest prime-factor counts ($\omega(n)=8{,}138$
forces $\ge 51{,}798$ digits and $P(n) > 10^{3{,}685}$; $\omega(n)=8{,}146$
forces $\ge 41{,}895$ digits and $P(n) > 10^{103}$; $\omega(n)=8{,}154$ forces
$\ge 41{,}128$ digits). Finally, we prove a sharp structural barrier for
two-giant-prime completions $n = M_0 P_1 P_2$ ($N_0 \ge 3$ odd and
$P_1 P_2 < t_0(N_0)^3 M_0^2 + t_0(N_0)^2 M_0 < 0.4523 M_0^2$ for
$M_0 \ge 10^5$), explaining why local CRT methods alone cannot settle the full
conjecture. All computational claims are certified by a standalone zero-float
Python standard-library verifier checking $846$ exact lattice witnesses in
under $16$ seconds; no interactive proof assistant was used.
\end{abstract}

\section{Introduction}\label{sec:intro}

In 1950, Giuga~\cite{Giuga1950} conjectured that an integer $n > 1$ is prime if
and only if
\begin{equation}\label{eq:giuga-power}
  \sum_{j=1}^{n-1} j^{n-1} \equiv -1 \pmod n.
\end{equation}
Agoh~\cite{Agoh1995} independently formulated the equivalent Bernoulli-number
congruence $n B_{n-1} \equiv -1 \pmod n$, and the assertion that no composite
$n$ satisfies~\eqref{eq:giuga-power} is now known as the \emph{Agoh--Giuga
conjecture}. By Giuga's characterization~\cite{Giuga1950} and Korselt's
criterion for Carmichael numbers~\cite{Korselt1899} (see also Borwein et
al.~\cite{Borwein_1996} and Luca, Pomerance, and
Shparlinski~\cite{LucaPomeranceShparlinski2009}), a composite integer $n$
satisfies~\eqref{eq:giuga-power} if and only if $n$ is odd, squarefree, and
satisfies, for every prime $p \mid n$,
\begin{equation}\label{eq:local}
  p \mid \frac{n}{p} - 1 \quad\text{(Giuga condition)}
  \qquad\text{and}\qquad
  (p - 1) \mid (n - 1) \quad\text{(Carmichael condition)}.
\end{equation}
Equivalently, letting $S = \PP(n) = \{p \text{ prime} : p \mid n\}$ and
$m = \omega(n) = |S|$, the Giuga condition is equivalent to the reciprocal
identity
\begin{equation}\label{eq:recip}
  \sum_{p \in S} \frac{1}{p} - \frac{1}{n} = a
\end{equation}
for a positive integer $a \ge 1$ (called the \emph{value} of $n$), while the
Carmichael condition implies that $S$ is a \emph{normal} set of odd primes: no
two distinct $p, q \in S$ satisfy $p \mid (q - 1)$. Indeed, $p \mid (q - 1)$
and $(q - 1) \mid (n - 1)$ would give $p \mid (n - 1)$, contradicting $p \mid n$.

Lower bounds on a hypothetical counterexample have a long computational history:
Giuga~\cite{Giuga1950} obtained $1{,}000$ digits, Bedocchi~\cite{Bedocchi1985}
reached $1{,}700$ digits, Borwein, Borwein, Borwein, and
Girgensohn~\cite{Borwein_1996} pushed normal-set tree pruning to depth $135$
($\omega(n) \ge 3{,}459$, $>10^{13{,}886}$), and Borwein, Maitland, and
Skerritt~\cite{DBLP:journals/int/BorweinMS13} reached depth $311$ in over ten
days of computation and $120\,\mathrm{GB}$ of storage, proving
$\omega(n) \ge 4{,}771$ and $n > 10^{19{,}907}$. Crucially, Borwein et
al.~\cite{Borwein_1996,DBLP:journals/int/BorweinMS13} observed that the greedy
normal set $L = L_{27,692}$ starting from $\{5,7\}$ up to the $27{,}692$nd odd
prime $q_{27,692} = 321{,}053$ already has $|L| = 8{,}135$ and
$\sum_{p\in L} 1/p > 1$, creating an absolute \emph{normal-set ceiling} of
$8{,}135$ prime factors for pure normality arguments. Because their trimmed
search tree grew exponentially with depth, they remarked that exhausting their
algorithm up to $d = 27{,}692$ appeared ``exceptionally unlikely to
happen''~\cite[p.~7]{DBLP:journals/int/BorweinMS13}. Recent work through
October 2026 continues to investigate modular and algebraic obstructions without
resolving the conjecture: Mavrotheris~\cite{Mavrotheris2026} derived mod-$8$
congruences via unimodular star lattices and van der Blij's theorem, and
L\'opez~\cite{Lopez2026} formulated an idempotence criterion for
counterexamples.

\subsection*{Summary of results and honest scope}

The full Agoh--Giuga conjecture remains open; our results are partial theorems
that (i) completely exhaust the normal-set extremal problem with a short
elementary proof and a $7$-leaf partition, and (ii) break decisively past the
normal-set ceiling by incorporating the exact Giuga and Carmichael congruences
via forced-prime Chinese Remainder Theorem (CRT) moduli and $2$-dimensional
lattice refutation. Specifically, within the scope of literature indexed through
October 2026, we establish the following new results:

\begin{enumerate}
  \item \textbf{Exact Extremal Normal-Set Theorem (Section~\ref{sec:starpack}).}
  Using a four-line \emph{Star-Pack Reciprocal Dominance Lemma}
  (Lemma~\ref{lem:starpack}) on a $7$-leaf partition of the odd primes, we prove
  (Theorem~\ref{thm:extremal}) that every finite normal set $S \subset \Podd$
  with $\sum_{p\in S} 1/p > 1$ satisfies $|S| \ge 8{,}135$ (not uniquely
  attained), $\max(S) \ge 321{,}053$ (equality iff $S = L$), and
  $\prod_{p\in S} p \ge P(L)$ ($40{,}733$ digits, equality iff $S = L$). This
  exhausts the Borwein--Maitland--Skerritt algorithm ($j_d = 8{,}135$ for all
  $d \ge 27{,}692$) and resolves the extremal lower-bound consequence of Borwein
  et al.~\cite[Open Problem~5]{Borwein_1996} (without settling the intermediate
  relaxed tree scores $j_d = j_d(L_d)$ for $312 \le d \le 27{,}691$). Moreover,
  outside the extremal leaf $\{5,7,23,83\}\subseteq S$, $3\notin S$, we prove
  $|S| \ge 8{,}846$, $\prod_{p\in S} p \ge 10^{44{,}627}$, and the exact
  attained bound $\max(S) \ge 352{,}757$.

  \item \textbf{Elementary mod-$16$ congruence and $a \ge 2$ splitting
  (Section~\ref{sec:elementary}).} Combining quadratic reciprocity with an
  elementary $+8$-invariance calculation in $(\mathbb{Z}/16\mathbb{Z})^\times$,
  we give a short self-contained proof of the mod-$8$ congruence
  $m - a \equiv 1 \pmod 8$ (first proved for odd Giuga numbers by
  Mavrotheris~\cite{Mavrotheris2026}) together with its exact mod-$16$
  refinement (Theorem~\ref{thm:mod16}). For $a \ge 2$, a purely elementary
  three-line splitting argument (Theorem~\ref{thm:a-ge-2}) using zero analytic
  number theory yields $\omega(n) \ge 329{,}188$, $P(n) \ge 1{,}284{,}211$, and
  at least $1{,}808{,}611$ decimal digits.

  \item \textbf{Unconditional Tier 2 exclusion bounds and dichotomies
  (Section~\ref{sec:tier2}).} Applying forced-prime CRT congruences
  $X \equiv C_F Y \pmod{M_F}$, sound suffix-minimum $2$-dimensional integer
  knapsacks, linear forcing cascades, zero-float Lagrangian star inequalities,
  and $846$ certified $2$-dimensional lattice witnesses, we prove
  (Theorem~\ref{thm:tier2}) that every composite Agoh--Giuga counterexample $n$
  satisfies
  \[
    \omega(n) \ge 8{,}138,
    \qquad
    P(n) \ge 326{,}057,
    \qquad
    n \ge 10^{40{,}921} \text{ ($\ge 40{,}922$ decimal digits)},
  \]
  and that small factor counts force massive prime factors and digit lengths
  ($\omega(n)=8{,}138 \implies \ge 51{,}798$ digits and $P(n)>10^{3{,}685}$;
  $\omega(n)=8{,}146 \implies \ge 41{,}895$ digits and $P(n)>10^{103}$;
  $\omega(n)=8{,}154 \implies \ge 41{,}128$ digits).

  \item \textbf{Sharp $r=2$ giant-prime structural barrier
  (Section~\ref{sec:barrier}).} For any $a = 1$ completion $n = M_0 P_1 P_2$
  with $N_0 = M_0(1 - \sum_{p\mid M_0} 1/p)$, we prove
  (Theorem~\ref{thm:barrier}) that $N_0 \ge 3$ is odd and
  $P_1 P_2 < t_0(N_0)^3 M_0^2 + t_0(N_0)^2 M_0$, where
  $t_0(N_0) = \bigl(1 + \sqrt{1+4N_0}\bigr)/(2N_0)$ (giving
  $P_1 P_2 < 0.4523 M_0^2$ whenever $M_0 \ge 10^5$). Since
  $M_L = P(L)\lambda(L)$ has $54{,}909$ digits while $M_0^2 \approx 10^{81{,}464}$,
  this pinpoints why local CRT + lattice methods leave
  $3 \le N_0 \lesssim 10^{17{,}704}$ open.
\end{enumerate}

To separate discovery from certification cleanly, we distinguish theoretical
lemmas proved purely algebraically from computationally certified theorems
checked by a standalone, zero-float, fail-closed Python verifier
(\texttt{src/verify\_agoh\_giuga\_certificate.py}) consuming a serialized exact
certificate (\texttt{certificates/agoh\_giuga\_tier2\_certificate.json}).

\section{Star-Pack Dominance and the Extremal Normal-Set Theorem}
\label{sec:starpack}

Let $q_1 = 3 < q_2 = 5 < q_3 = 7 < \cdots$ denote the odd primes, so that
$q_{27,692} = 321{,}053$. Define $L = L_{27,692}$ to be the greedy normal set
constructed from $\{5,7\}$ up to $B := 321{,}053$: initialize $L^{(0)} = \{5,7\}$
and scan odd primes $p \in [11, B]$ in ascending order, adjoining $p$ whenever
no already chosen prime $r < p$ divides $p - 1$. (In Borwein, Maitland, and
Skerritt~\cite[Definition~11]{DBLP:journals/int/BorweinMS13}, the index is
shifted by one via $Q_k = \{q_1,\dots,q_{k-1}\}$, though the text likewise
refers to this $8{,}135$-element set as $L_{27,692}$.) Direct exact rational
evaluation shows $|L| = 8{,}135$, $\max(L) = 321{,}053$,
\[
  \sum_{p \in L} \frac{1}{p} \approx 1.000000918503 > 1,
  \qquad
  \sum_{p \in L \setminus \{321{,}053\}} \frac{1}{p} \approx 0.999997803752 < 1,
\]
and $\gap_L := \sum_{p\in L} 1/p - 1 < 1/B$. Its product $P(L) := \prod_{p\in L} p$
satisfies $10^{40{,}732} \le P(L) < 10^{40{,}733}$ ($40{,}733$ decimal digits,
$135{,}311$ bits).

The key to replacing exponential depth-$27{,}692$ tree pruning by a $7$-leaf
exact proof is a four-line graph-theoretic upper bound on normal subsets of an
interval $[1,B]$.

\begin{lemma}[Star-Pack Reciprocal Dominance]\label{lem:starpack}
Let $N, F \subset \Podd$ be disjoint finite sets with $N$ normal, let
$B \ge \max(N \cup F \cup \{3\})$, and set
\[
  U := \bigl\{p \in \Podd \cap [1, B] \setminus (N \cup F) :
  \forall r \in N,\; r \nmid (p - 1) \text{ and } p \nmid (r - 1)\bigr\}.
\]
Suppose $U = \bigsqcup_{c\in C} A_c$ is partitioned into pairwise disjoint
\emph{stars} $A_c = \{c\} \sqcup T_c$ such that for every $c \in C$:
\begin{enumerate}
  \item[(i)] $c \mid (q - 1)$ (hence $q \ge 2c + 1 > c$) for every $q \in T_c$;
  \item[(ii)] $\capf(c) := \frac{1}{c} - \sum_{q\in T_c} \frac{1}{q} > 0$.
\end{enumerate}
Let $\Sigma := \sum_{p \in N \cup C} 1/p$. Let $S \subset \Podd$ be any finite
set with $\sum_{p\in S} 1/p > 1$, $N \subseteq S$, $S \cap F = \emptyset$, and
such that $S_{\le B} := S \cap [1, B]$ is normal (no normality is assumed on
$S_{>B} := S \cap (B, \infty)$). Partition $C = M_0 \sqcup M_1 \sqcup M_+$ by
\[
  M_0 := \{c \in C : S \cap A_c = \emptyset\},
  \quad
  M_1 := \{c \in C : \emptyset \neq S \cap A_c \subseteq T_c\},
  \quad
  M_+ := \{c \in C : c \in S\}.
\]
Then:
\begin{enumerate}
  \item For every $c \in C$, $\sum_{q \in S \cap A_c} 1/q \le 1/c$, with strict
  inequality for $c \in M_0 \cup M_1$; in particular
  $\sum_{q \in S_{\le B}} 1/q \le \Sigma$.
  \item If $\Sigma \le 1$, then $S_{>B} \neq \emptyset$, $\max(S) > B$, and with
  $t := \lfloor B(1 - \Sigma)\rfloor + 1 \ge 1$ we have
  \[
    |S_{>B}| \ge |M_0| + t,
    \qquad
    |S| \ge |N \cup C| + t,
    \qquad
    \prod_{q\in S} q > \left(\prod_{p \in N \cup C} p\right) B^t.
  \]
  \item If $\Sigma > 1$ and
  $\Sigma - 1 < \min\bigl(\frac{1}{B},\, \min_{c\in C} \capf(c)\bigr)$, then
  $|S_{>B}| \ge |M_0|$, $|S| \ge |N \cup C|$, $\max(S) \ge B$ (with
  $\max(S) \le B$ iff $S = N \cup C$), and $\prod_{q\in S} q \ge \prod_{p\in N\cup C} p$
  (with equality iff $S = N \cup C$).
\end{enumerate}
\end{lemma}

\begin{proof}
Since $S_{\le B}$ is normal and $N \subseteq S_{\le B}$, we have
$S_{\le B} \subseteq N \sqcup U$. Within each star $A_c = \{c\}\sqcup T_c$,
condition~(i) forbids $S_{\le B}$ from containing both $c$ and any element of
$T_c$, so $M_0, M_1, M_+$ partition $C$ and condition~(ii) yields
$\sum_{q\in S\cap A_c} 1/q \le 1/c$, proving~(1). Using $1/c \ge 1/B$ for
$c \in M_0$ and $1/q < 1/B$ for $q \in S_{>B}$, we obtain
\begin{equation}\label{eq:star-balance}
  1 < \sum_{q\in S}\frac{1}{q}
  \le \Sigma - \sum_{c\in M_1}\capf(c) - \sum_{c\in M_0}\frac{1}{c}
      + \sum_{q\in S_{>B}}\frac{1}{q}
  \le \Sigma - \frac{|M_0|}{B} + \frac{|S_{>B}|}{B},
\end{equation}
where the last inequality is strict whenever $S_{>B} \neq \emptyset$.

In~(2), $\Sigma \le 1$ forces $S_{>B} \neq \emptyset$, hence $\max(S) > B$ and
$|S_{>B}| > |M_0| + B(1 - \Sigma)$, so $|S_{>B}| \ge |M_0| + t$. Since each star
in $M_+ \cup M_1$ contributes at least one prime to $S_{\le B}$,
$|S_{\le B}| \ge |N \cup C| - |M_0|$, whence $|S| = |S_{\le B}| + |S_{>B}|
\ge |N \cup C| + t$. For the product, each $c \in M_1$ satisfies
$\prod_{q \in S\cap T_c} q \ge 2c + 1 > c$, while
$\prod_{q\in S_{>B}} q > B^{|M_0|+t} \ge (\prod_{c\in M_0} c) B^t$; multiplying
over $N, M_+, M_1$, and $S_{>B}$ gives the product inequality.

In~(3), suppose first that $S_{>B} = \emptyset$. If $S \neq N \cup C$, then
$M_0 \cup M_1 \neq \emptyset$, so~\eqref{eq:star-balance} drops by at least
$\min(1/B, \min_{c\in C}\capf(c)) > \Sigma - 1$, contradicting
$\sum_{q\in S} 1/q > 1$. Hence $S_{>B} = \emptyset$ forces $S = N \cup C$,
giving the $\max(S)$ statement. If $S_{>B} \neq \emptyset$,
then~\eqref{eq:star-balance} yields $|S_{>B}| > |M_0| - B(\Sigma - 1) > |M_0| - 1$,
so $|S_{>B}| \ge |M_0|$. Therefore $|S| \ge |N \cup C|$, and since
$S_{>B} \neq \emptyset$ implies $\prod_{q\in S_{>B}} q > B^{|M_0|} \ge \prod_{c\in M_0} c$
strictly, we obtain $\prod_{q\in S} q > \prod_{p\in N\cup C} p$.
\end{proof}

Note that Lemma~\ref{lem:starpack} does \emph{not} require the center set $C$
itself to be normal for parts~(1)--(2) to hold: any valid disjoint star
partition of $U$ gives a rigorous upper bound. To construct a star partition
mechanically for any leaf $(N,F)$ and cutoff $B$, we use \emph{max-capacity
ascending Star-Pack}: scan $p \in U$ in ascending order; if any already chosen
center $c < p$ divides $p - 1$ and has remaining capacity strictly exceeding
$1/p$, assign $p$ to the leaf set $T_c$ of the divisor $c$ with largest
remaining capacity; otherwise make $p$ a new center with $T_p = \emptyset$.
Whenever no center divides $p - 1$ for any other center (zero non-normal
centers), the resulting set $N \cup C$ coincides identically with the greedy
normal set extending $N$ within $[1,B]\setminus F$.

\begin{theorem}[Exact Extremal Normal-Set Theorem; computationally certified]
\label{thm:extremal}
Partition all subsets of $\Podd$ into the seven mutually exclusive, exhaustive
leaves of Table~\ref{tab:leaves} according to membership in $\{3,5,7,23,83\}$.
Then:
\begin{enumerate}
  \item Every finite normal set $S \subset \Podd$ with $\sum_{p\in S} 1/p > 1$
  satisfies
  \[
    |S| \ge 8{,}135,
    \qquad
    \max(S) \ge 321{,}053,
    \qquad
    \prod_{p\in S} p \ge P(L) > 10^{40{,}732}.
  \]
  Equality holds in $\max(S) = 321{,}053$ if and only if $S = L$, and in
  $\prod_{p\in S} p = P(L)$ if and only if $S = L$. Equality in $|S| = 8{,}135$
  is \emph{not} uniquely attained: the normal set
  $L' = (L \setminus \{321{,}053\}) \cup \{321{,}073\}$ also satisfies
  $|L'| = 8{,}135$ and $\sum_{p\in L'} 1/p > 1$.
  \item \emph{(Exhaustion of Borwein--Maitland--Skerritt~\cite{DBLP:journals/int/BorweinMS13}.)}
  In the notation of~\cite[Definition~7]{DBLP:journals/int/BorweinMS13},
  $j_d = 8{,}135$ for every $d \ge 27{,}692$.
  \item \emph{(Strict separation outside Leaf~2.)} Every finite normal set $S$
  with $\sum_{p\in S} 1/p > 1$ outside Leaf~2 (i.e., not satisfying
  $\{5,7,23,83\}\subseteq S$ and $3\notin S$) obeys
  \[
    |S| \ge 8{,}846,
    \qquad
    \prod_{p\in S} p \ge 10^{44{,}627} \text{ ($\ge 44{,}628$ digits)},
    \qquad
    \max(S) \ge 352{,}757.
  \]
  The bound $\max(S) \ge 352{,}757$ is sharp and uniquely attained outside
  Leaf~2 by the $8{,}880$-element greedy normal set $L^{(3)}$ of Leaf~3
  ($N=\{5,7,23\}$, $F=\{3,83\}$).
  \item Max-capacity ascending Star-Pack on Leaf~2 has zero non-normal centers
  throughout $[1, 352{,}756]$ ($|L_{338,032}| = 8{,}512$,
  $|L_{352,756}| = 8{,}855$).
\end{enumerate}
\end{theorem}

\begin{table}[t]
\centering
\caption{Exact $7$-leaf max-capacity ascending Star-Pack partition at $B = 321{,}053$
(all bounds certified in exact rational/integer arithmetic). Outside Leaf~2, every
leaf also satisfies $\Sigma < 1$ at $B_{\mathrm{out}} = 352{,}756$, while Leaf~3
attains normality and $\Sigma > 1$ uniquely at $352{,}757$.}
\label{tab:leaves}
\small
\begin{tabular}{cllrrrrr}
\toprule
Leaf & Forced $N$ & Forbidden $F$ & $|N\cup C|$ & Non-norm. & $t$ & Min $|S|$ & Min digits \\
\midrule
1 & $\{3\}$ & $\emptyset$ & $8{,}009$ & $2{,}299$ & $6{,}748$ & $14{,}757$ & $77{,}339$ \\
2 & $\{5,7,23,83\}$ & $\{3\}$ & $8{,}135$ & $0$ & $0$ & $8{,}135$ & $40{,}733$ \\
3 & $\{5,7,23\}$ & $\{3,83\}$ & $8{,}156$ & $0$ & $690$ & $8{,}846$ & $44{,}628$ \\
4 & $\{5,7\}$ & $\{3,23\}$ & $8{,}286$ & $17$ & $4{,}355$ & $12{,}641$ & $65{,}487$ \\
5 & $\{7\}$ & $\{3,5\}$ & $10{,}455$ & $1{,}882$ & $3{,}239$ & $13{,}694$ & $70{,}327$ \\
6 & $\{5\}$ & $\{3,7\}$ & $10{,}174$ & $1{,}453$ & $16{,}000$ & $26{,}174$ & $139{,}312$ \\
7 & $\emptyset$ & $\{3,5,7\}$ & $12{,}304$ & $3{,}347$ & $12{,}116$ & $24{,}420$ & $128{,}610$ \\
\bottomrule
\end{tabular}
\end{table}

\begin{proof}
Exhaustiveness and mutual exclusivity of the $7$ leaves are verified directly
over all $2^5 = 32$ truth assignments to $\{3,5,7,23,83\}$. Running max-capacity
ascending Star-Pack at $B = 321{,}053$ and independently checking every star's
exact rational sum $\sum_{q\in T_c} 1/q < 1/c$ yields Table~\ref{tab:leaves}. On
Leaf~2, the star packing has zero non-normal centers, so $N \cup C = L$ with
$\Sigma = \sum_{p\in L} 1/p > 1$, $\Sigma - 1 < 1/B$, and
$\min_{c\in C}\capf(c) = 1/B$ (attained at $c = B$, which has $T_B = \emptyset$).
Applying Lemma~\ref{lem:starpack}(3) on Leaf~2 and Lemma~\ref{lem:starpack}(2)
on Leaves~1 and~3--7 proves part~(1), together with exact verification of the
alternative $8{,}135$-element normal set $L'$.

For part~(2), recall that in~\cite{DBLP:journals/int/BorweinMS13} a vertex
$N \in \mathcal{N}_{27,692}$ is a normal subset of $Q_{27,692} = \{3,\dots,321{,}037\}$
and $T_{27,692}(N)$ appends odd primes $\ge q_{27,692} = B$ compatible with $N$
until the reciprocal sum exceeds $1$. Every appended prime satisfies $1/q \le 1/B$,
and $T_{27,692}(N) \cap [1, B - 1] = N$ is normal. Because
Lemma~\ref{lem:starpack} never assumes mutual normality among primes $\ge B$
(and the non-strict bound $1/q \le 1/B$ still gives
$|T_{27,692}(N)| > |L| - B(\Sigma - 1) > 8{,}134$ on Leaf~2 and $>8{,}845$ on
Leaves~1, 3--7), we have $j_{27,692}(N) \ge 8{,}135$ for all
$N \in \mathcal{N}_{27,692}$. Combining this with monotonicity
\cite[Theorem~10]{DBLP:journals/int/BorweinMS13} and the explicit normal set $L$
gives $8{,}135 \le j_{27,692} \le j_d \le |L| = 8{,}135$ for all $d \ge 27{,}692$.

For part~(3), running max-capacity ascending Star-Pack at $B_{\mathrm{out}} = 352{,}756$
on Leaves~1 and~3--7 certifies $\Sigma < 1$ on every non-Leaf-2 leaf, so
Lemma~\ref{lem:starpack}(2) gives $\max(S) \ge 352{,}757$ outside Leaf~2. At
$B = 352{,}757$, Leaf~3 produces zero non-normal centers, $8{,}880$ elements,
$\Sigma > 1$, and $\min_{c\in C}\capf(c) > \Sigma - 1$; hence
Lemma~\ref{lem:starpack}(3) proves unique attainment at $L^{(3)}$. Part~(4) is
checked directly at $B_{\mathrm{out}} = 352{,}756$.
\end{proof}

\begin{remark}\label{rem:open5}
Theorem~\ref{thm:extremal}(2) establishes the full lower-bound consequence of
Open Problem~5 of Borwein et al.~\cite{Borwein_1996} (and sharpens their
$36{,}069$-digit estimate to the exact normal-set minimum of $40{,}733$ digits),
while leaving open whether $j_d = j_d(L_d)$ holds at every intermediate depth
$312 \le d \le 27{,}691$.
\end{remark}

\section{Elementary Mod-16 Congruences and Value Splitting (\texorpdfstring{$a \ge 2$}{a >= 2})}
\label{sec:elementary}

To break past the $8{,}135$-prime normal-set ceiling, we exploit the Giuga
equation~\eqref{eq:recip}. Let $n = \prod_{p\in S} p$ be an odd weak Giuga
number ($p \mid n/p - 1$ for all $p \in S$, no Carmichael condition required),
with $m = |S|$ and value $a \in \mathbb{Z}_{\ge 1}$. Let
$c_k := \#\{p \in S : p \equiv k \pmod 8\}$ for $k \in \{1,3,5,7\}$, and let
$r := c_3 + c_7 = \#\{p \in S : p \equiv 3 \pmod 4\}$.

\begin{theorem}[Mod-$16$ and Mod-$8$ Congruences]\label{thm:mod16}
For every odd squarefree integer $n = \prod_{p\in S} p$ satisfying
$\sum_{p\in S} 1/p - 1/n = a \in \mathbb{Z}_{\ge 1}$:
\begin{enumerate}
  \item $r \equiv 0 \text{ or } 1 \pmod 4$;
  \item modulo $16$,
  \[
    m - a - 1 \equiv
    \begin{cases}
      2r + 8c_3 \pmod{16} & \text{if } r \equiv 0 \pmod 4,\\
      2(r - 1) + 8c_5 \pmod{16} & \text{if } r \equiv 1 \pmod 4;
    \end{cases}
  \]
  \item in particular, $m - a \equiv 1 \pmod 8$ (cf.\ Mavrotheris~\cite{Mavrotheris2026}).
\end{enumerate}
\end{theorem}

\begin{proof}
For part~(1), the Giuga condition gives $\prod_{q\in S\setminus\{p\}} q \equiv 1 \pmod p$
for each $p \in S$. Taking Legendre symbols and multiplying over all $p \in S$,
quadratic reciprocity yields
\[
  1 = \prod_{p\in S} \prod_{q\neq p}\left(\frac{q}{p}\right)
    = \prod_{\{p,q\}\subset S}\left(\frac{p}{q}\right)\left(\frac{q}{p}\right)
    = (-1)^{\binom{r}{2}},
\]
so $\binom{r}{2}$ is even, i.e., $r \equiv 0$ or $1 \pmod 4$.

For part~(2), reducing~\eqref{eq:recip} in $(\mathbb{Z}/16\mathbb{Z})^\times$
gives $a \equiv \sum_{p\in S} p^{-1} - n^{-1} \pmod{16}$. Observe that for any
odd $p$, $(p + 8)(p^{-1} + 8) = 1 + 8(p + p^{-1}) + 64 \equiv 1 \pmod{16}$ since
$p + p^{-1}$ is even. Thus shifting any factor $p \mapsto p + 8$ shifts both
$p^{-1}$ and $n^{-1} = (n/p)^{-1}p^{-1}$ by $+8 \pmod{16}$, leaving
$\sum_{p\in S} p^{-1} - n^{-1} \pmod{16}$ invariant. We may therefore replace
each prime by its representative in $\{1,3,5,7\} \pmod 8$. Using
$1^{-1}\equiv 1$, $3^{-1}\equiv 11$, $5^{-1}\equiv 13$, $7^{-1}\equiv 7 \pmod{16}$,
we have
\[
  \sum_{p\in S} p^{-1} \equiv c_1 + 11c_3 + 13c_5 + 7c_7
  = m + 10c_3 + 12c_5 + 6c_7 \pmod{16}.
\]
In $(\mathbb{Z}/16\mathbb{Z})^\times$, we have $11 \equiv 7\cdot 13$, $7^2\equiv 1$,
and $13^k = (1 - 4)^k \equiv 1 - 4k \pmod{16}$ (since $4^2 = 16 \equiv 0 \pmod{16}$).
Hence
\[
  n^{-1} \equiv 11^{c_3} 13^{c_5} 7^{c_7}
  \equiv 7^{c_3 + c_7} 13^{c_3 + c_5}
  \equiv 7^r\bigl(1 - 4(c_3 + c_5)\bigr) \pmod{16}.
\]
If $r \equiv 0 \pmod 4$, then $7^r \equiv 1 \pmod{16}$, so
\[
  a \equiv (m + 10c_3 + 12c_5 + 6c_7) - (1 - 4c_3 - 4c_5)
    \equiv m - 1 + 6r + 8c_3 \pmod{16}.
\]
Since $4r \equiv 0 \pmod{16}$, we have $-6r \equiv 2r \pmod{16}$ and
$-8c_3 \equiv 8c_3 \pmod{16}$, giving $m - a - 1 \equiv 2r + 8c_3 \pmod{16}$.
If $r \equiv 1 \pmod 4$, then $7^r \equiv 7 \pmod{16}$ and
$7(1 - 4(c_3+c_5)) \equiv 7 + 4(c_3+c_5) \pmod{16}$, giving
\[
  a \equiv m - 7 + 6c_3 + 8c_5 + 6c_7 = m - 1 + 6(r - 1) + 8c_5 \pmod{16},
\]
and hence $m - a - 1 \equiv 2(r - 1) + 8c_5 \pmod{16}$. Reducing modulo $8$
gives~(3) immediately.
\end{proof}

\begin{remark}
Mavrotheris~\cite{Mavrotheris2026} derived the odd-Giuga congruence
$m - a \equiv 1 \pmod 8$ using unimodular star lattices and van der Blij's
characteristic-vector theorem; Theorem~\ref{thm:mod16} provides an elementary
$10$-line proof and lifts the identity modulo~$16$. Note that when $a = 1$,
Theorem~\ref{thm:mod16}(3) forces $m \equiv 2 \pmod 8$, lifting
Theorem~\ref{thm:extremal}'s $|S| \ge 8{,}135$ immediately to
$\omega(n) \ge 8{,}138$.
\end{remark}

We next separate the value $a \ge 2$ regime by a purely elementary argument
requiring zero analytic prime-counting estimates.

\begin{theorem}[Elementary $a \ge 2$ Splitting Bound]\label{thm:a-ge-2}
Every composite Agoh--Giuga counterexample $n$ of value $a \ge 2$ satisfies
\[
  \omega(n) \ge 329{,}188,
  \qquad
  P(n) \ge 1{,}284{,}211,
  \qquad
  n > 10^{1{,}808{,}610} \text{ ($\ge 1{,}808{,}611$ decimal digits)}.
\]
\end{theorem}

\begin{proof}
Write $S = \PP(n) = \{p_1 < p_2 < \dots < p_m\}$. Since $\sum_{i=1}^m 1/p_i > a \ge 2$,
let $k$ be the least index such that $\sum_{i=1}^k 1/p_i > 1$, and set
$S_1 = \{p_1,\dots,p_k\}$ and $S_2 = \{p_{k+1},\dots,p_m\}$. Because $S_1$ is
normal with $\sum_{p\in S_1} 1/p > 1$, Theorem~\ref{thm:extremal}(1) gives
$k \ge 8{,}135$, $p_k \ge 321{,}053$, and $\prod_{p\in S_1} p \ge P(L) \ge 10^{40{,}732}$.
By minimality of $k$, $\sum_{i=1}^{k-1} 1/p_i \le 1$, so
$\sum_{p\in S_1} 1/p \le 1 + 1/p_k$. Hence
\[
  1 - \frac{1}{p_k} < \left(a + \frac{1}{n}\right) - \sum_{p\in S_1}\frac{1}{p}
  = \sum_{p\in S_2}\frac{1}{p} < \frac{|S_2|}{p_k},
\]
forcing $|S_2| > p_k - 1$, i.e., $|S_2| \ge p_k \ge 321{,}053$. Thus
$\omega(n) = k + |S_2| \ge 8{,}135 + 321{,}053 = 329{,}188$.

Every prime $p \in S_2$ exceeds $321{,}053$ and is coprime to $6$. Since
$321{,}053 = 6\cdot 53{,}508 + 5$, the integers $> 321{,}053$ coprime to $6$
alternate by $+2$ and $+4$, so the $321{,}053$rd such integer is
$321{,}053 + 6\cdot 160{,}526 + 2 = 1{,}284{,}211$ (which is prime), giving
$P(n) \ge 1{,}284{,}211$. Finally, exact integer exponentiation confirms
$321{,}053^{10{,}000} > 10^{55{,}065}$ and $321{,}053^{1{,}053} > 10^{5{,}798}$,
whence
\begin{align*}
  n &> P(L)\cdot 321{,}053^{321{,}053}
     = P(L)\cdot\bigl(321{,}053^{10{,}000}\bigr)^{32}\cdot 321{,}053^{1{,}053} \\
    &> 10^{40{,}732 + 32\cdot 55{,}065 + 5{,}798} = 10^{1{,}808{,}610}. \qedhere
\end{align*}
\end{proof}

\section{Breaking Past the Normal-Set Ceiling via Forced-Prime CRT}
\label{sec:tier2}

Theorems~\ref{thm:extremal} and~\ref{thm:a-ge-2} reduce all unconditional
exclusion bounds below $8{,}846$ prime factors, largest prime $352{,}757$, or
$44{,}628$ decimal digits to the single case $a = 1$ inside Leaf~2
($\{5,7,23,83\}\subseteq S$, $3\notin S$). On this leaf, any normal base set
$L_{\mathrm{base}} = L_{P_0}$ ($P_0 \le 352{,}756$) turns the simultaneous
congruences~\eqref{eq:local} into a massive $2$-dimensional lattice constraint.

\begin{lemma}[Forced-Prime CRT and $2$D Lattice Box Refutation]\label{lem:crt-lattice}
Let $L_{\mathrm{base}} \subset \Podd$ be a finite normal set with
$P_{\mathrm{base}} := \prod_{p\in L_{\mathrm{base}}} p$. For any composite
Agoh--Giuga counterexample $n = \prod_{p\in S} p$, define the coprime positive
integers
\[
  Y := \prod_{p \in L_{\mathrm{base}} \setminus S} p,
  \qquad
  X := \prod_{q \in S \setminus L_{\mathrm{base}}} q,
  \qquad
  \text{so that } Y n = P_{\mathrm{base}} X \text{ and } \gcd(X, Y) = 1.
\]
Let $F \subseteq S \cap L_{\mathrm{base}}$ be any nonempty subset of forced
primes, and set $\lambda(F) := \operatorname{lcm}_{p\in F}(p - 1)$ and
$M_F := \bigl(\prod_{p\in F} p\bigr)\lambda(F)$. Then:
\begin{enumerate}
  \item Normality of $L_{\mathrm{base}}$ ensures that the moduli $\{p : p\in F\}$
  and $\lambda(F)$ are pairwise coprime and $\gcd(P_{\mathrm{base}}, \lambda(F)) = 1$.
  Consequently, the Giuga and Carmichael congruences~\eqref{eq:local} at $p\in F$
  combine via CRT into a unique residue $C_F \in (\mathbb{Z}/M_F\mathbb{Z})^\times$
  such that
  \begin{equation}\label{eq:crt-main}
    X \equiv C_F Y \pmod{M_F}.
  \end{equation}
  \item Let $(u, v) \in \mathbb{Z}^2 \setminus \{(0,0)\}$ satisfy
  $u \equiv C_F v \pmod{M_F}$, and let $U, R \ge 1$ be integers. If either
  \begin{enumerate}
    \item[(a)] $uv < 0$ and $|v|R + |u|U < M_F$, or
    \item[(b)] $uv > 0$, $|v|R < M_F$, $|u|U < M_F$, and
    $\max\bigl(|u|/g - R,\, |v|/g - U\bigr) > 0$ where $g = \gcd(|u|, |v|)$,
  \end{enumerate}
  then no coprime integer pair $(X, Y) \in [1, R] \times [1, U]$ satisfies~\eqref{eq:crt-main}.
\end{enumerate}
\end{lemma}

\begin{proof}
For~(1), if a prime $q \in L_{\mathrm{base}}$ divided $p - 1$ for some
$p \in F \subseteq L_{\mathrm{base}}$, then $q \neq p$ would contradict normality
of $L_{\mathrm{base}}$. Thus $\gcd(P_{\mathrm{base}}, \lambda(F)) = 1$, and in
particular $\gcd(\prod_{p\in F} p, \lambda(F)) = 1$. For each $p \in F$, $p \nmid Y$
and $p^2 \nmid P_{\mathrm{base}}$, so the Giuga condition $n/p \equiv 1 \pmod p$
becomes $(P_{\mathrm{base}}/p) X \equiv Y \pmod p$, while the Carmichael
condition $n \equiv 1 \pmod{\lambda(F)}$ becomes
$P_{\mathrm{base}} X \equiv Y \pmod{\lambda(F)}$. Since every coefficient is
invertible modulo its modulus, CRT yields a unique $C_F \bmod M_F$ satisfying~\eqref{eq:crt-main}.

For~(2), $u \equiv C_F v$ and $X \equiv C_F Y \pmod{M_F}$ give
$vX - uY \equiv 0 \pmod{M_F}$. Under~(a), $vX - uY$ has constant sign and
satisfies $0 < |vX - uY| = |v|X + |u|Y \le |v|R + |u|U < M_F$, a contradiction.
Under~(b), $-M_F < -|u|U < vX - uY < |v|R < M_F$, forcing $vX = uY$; since
$\gcd(X,Y) = 1$, we must have $(X,Y) = (|u|/g, |v|/g)$, which lies outside
$[1,R]\times [1,U]$.
\end{proof}

Notice that checking a witness $(u,v)$ in Lemma~\ref{lem:crt-lattice}(2) takes
just one modular multiplication and two integer comparisons, with zero floating
point operations. In practice, whenever $4UR < M_F$, running the extended
Euclidean algorithm on $(M_F, C_F)$ down to the threshold
$r_i \le \lfloor \sqrt{M_F R / U}\rfloor$ produces a valid witness directly.

\begin{theorem}[Unconditional Tier 2 Bounds and Dichotomies; computationally certified]
\label{thm:tier2}
Every composite counterexample $n$ to the Agoh--Giuga conjecture satisfies
unconditionally
\[
  \omega(n) \ge 8{,}138,
  \qquad
  P(n) \ge 326{,}057,
  \qquad
  n \ge 10^{40{,}921} \text{ ($\ge 40{,}922$ decimal digits)}.
\]
Moreover, for the smallest admissible prime-factor counts ($a = 1$,
$\omega(n) \equiv 2 \pmod 8$):
\begin{enumerate}
  \item if $\omega(n) = 8{,}138$, then $n \ge 10^{51{,}797}$ ($\ge 51{,}798$
  digits) and $P(n) > 10^{3{,}685}$;
  \item if $\omega(n) = 8{,}146$, then $n \ge 10^{41{,}894}$ ($\ge 41{,}895$
  digits) and $P(n) > 10^{103}$;
  \item if $\omega(n) = 8{,}154$, then $n \ge 10^{41{,}127}$ ($\ge 41{,}128$
  digits).
\end{enumerate}
\end{theorem}

The proof splits cleanly into two independent exact certifications: one for
$P(n) \ge 326{,}057$ and one for $n \ge 10^{40{,}921}$ (together with the three
dichotomies).

\subsection{Certification of \texorpdfstring{$P(n) \ge 326{,}057$}{P(n) >= 326,057}}\label{sec:proof-maxp}

Suppose $n$ is a counterexample with $P(n) \le P_0 := 326{,}056$ (the next prime
is $326{,}057$). By Theorems~\ref{thm:extremal} and~\ref{thm:a-ge-2}, $a = 1$
and $S$ lies in Leaf~2. Let $L_0 = L_{P_0}$ be the Leaf-2 max-capacity Star-Pack
center set up to $P_0$: $|L_0| = 8{,}254$, $\max(L_0) = 325{,}999$, and $L_0$ is
$100\%$ normal by Theorem~\ref{thm:extremal}(4). Because $P(n) \le P_0$, we have
$S_{>P_0} = \emptyset$, so every omitted center $c \in M_0 \cup M_1$ draws
strictly from $\gap_0 := \sum_{p\in L_0} 1/p - 1 \approx 3.6865 \times 10^{-4}$:
\begin{equation}\label{eq:p0-budget}
  \sum_{c \in M_0} \frac{1}{c}
  + \sum_{c \in M_1} \left(\frac{1}{c} - \sum_{q\in S\cap T_c}\frac{1}{q}\right)
  < \gap_0.
\end{equation}
Fixing the exact integer scale $D = 10^{12}$, set
$B_{\mathrm{int}} := \lceil D\cdot\gap_0\rceil$ and
$\capf_{\mathrm{int}}(c) := \lfloor D\cdot\capf_{P_0}(c)\rfloor$. Any center
$c \in C_0$ with $\capf_{\mathrm{int}}(c) \ge B_{\mathrm{int}}$ cannot lie in
$M_0 \cup M_1$ and is therefore forced into $S$; together with $\{5,7,23,83\}$
this gives an initial forced set $F_0$ of size $104$.

For any state with forced set $F_{\mathrm{cur}} \subseteq S$ and optional
omitted prime $p_{\mathrm{omit}} \notin S$, we bound $Y = \prod_{c\in L_0\setminus S} c$
as a function of $k = |S \setminus L_0|$ via a sound $2$-dimensional integer
knapsack. Each remaining star $c$ with $T_c \neq \emptyset$ has cost at least
$\capf_{\mathrm{int}}(c)$ and contributes at most $|T_c|$ leaves to $k$; computing
$\mathrm{dp}[ds][k]$ (minimum total scaled cost using $ds$ stars with full leaf
counts summing to $k$, initialized with $p_{\mathrm{omit}}$ in $M_0$ or $M_1$)
up to the exact fractional-knapsack upper bound $K_{\mathrm{ub}}$, and then
taking the \emph{suffix minimum} $\mathrm{dp}_{\mathrm{sound}}[ds][k] :=
\min_{k' \ge k} \mathrm{dp}[ds][k']$, guarantees soundness for \emph{every}
partial leaf subset $W_c \subseteq T_c$ (since replacing $W_c$ by $T_c$ cannot
increase cost and only increases $k$). For each feasible $(ds,k)$, the remaining
budget $B_{\mathrm{int}} - \mathrm{dp}_{\mathrm{sound}}[ds][k]$ caps
$m_0 = |M_0 \setminus \{p_{\mathrm{omit}}\}|$, yielding an exact integer bound
$U(k) = \max_{ds} \bigl(p_{\mathrm{omit}} \prod_{j=1}^{ds} c^{\mathrm{star}}_{\mathrm{top},j}
\prod_{j=1}^{m_0} L_{0,\mathrm{top},j}\bigr)$.

We close $P_0 = 326{,}056$ by a \emph{linear forcing cascade}: iterate through
non-forced primes $p \in L_0 \setminus F_0$ in descending order of
$\capf_{\mathrm{int}}(p)$ (skipping only $p = 233$, which has $47$ leaves and
remains unconstrained). At each step, Branch~A assumes $p \notin S$: subtracting
$\capf_{\mathrm{int}}(p)$ immediately forces dozens to thousands of additional
primes $F_{\mathrm{ext}}(p) \subseteq S$, and $k$-bucketing $[0, k_{\max}]$
into contiguous intervals $[k_{\mathrm{lo}}, k_{\mathrm{hi}}]$ with
$U_{\mathrm{box}} = \max_{k\in [k_{\mathrm{lo}}, k_{\mathrm{hi}}]} U(k)$ and
$R_{\mathrm{box}} = P_0^{k_{\mathrm{hi}}}$ refutes Branch~A via
Lemma~\ref{lem:crt-lattice}. Branch~B then adjoins $p$ to $F_{\mathrm{cur}}$.
After $152$ Branch~A steps ($|F_{\mathrm{cur}}| = 256$), Branch~B is refuted
directly in $3$ buckets ($[0,84]$, $[85,102]$, $[103,104]$), totaling $153$
branches and $155$ exact $(U,R)$ witnesses.

\subsection{Certification of \texorpdfstring{$n \ge 10^{40{,}921}$}{n >= 10\textasciicircum 40921} and the \texorpdfstring{$\omega(n)$}{omega(n)} Dichotomies}
\label{sec:proof-digits}

By Theorems~\ref{thm:extremal}, \ref{thm:mod16}, and~\ref{thm:a-ge-2}, any
counterexample with fewer than $40{,}922$ decimal digits satisfies $a = 1$, lies
in Leaf~2, and has $|S| = 8{,}135 + \Delta m$ with $\Delta m \equiv 3 \pmod 8$,
$\Delta m \ge 3$. Work relative to $L_{\mathrm{base}} = L$ ($B = 321{,}053$).
Writing $m_0 = |M_0|$, $m_1 = |M_1|$, $E = \sum_{c\in M_1}(|S\cap T_c| - 1) \ge 0$,
and $s_{\mathrm{out}} = |S_{>B}| = m_0 + \Delta m - E$, rearranging the
Star-Pack identity~\eqref{eq:star-balance} yields two exact inequalities:
\begin{align}
  \sum_{c\in M_0}\left(\frac{1}{c} - \frac{1}{B}\right)
  + \sum_{c\in M_1} \delta(c)
  + \sum_{q\in S_{>B}}\left(\frac{1}{B} - \frac{1}{q}\right)
  &< \gap_L + \frac{\Delta m}{B},
  \label{eq:delta-budget}\\
  \sum_{c\in M_0}\frac{1}{c} + \sum_{c\in M_1}\capf(c)
  &< \gap_L + T_F[m_0 + \Delta m],
  \label{eq:tail-budget}
\end{align}
where $\delta(c) := \capf(c) + (|T_c| - 1)/B > 0$, and $T_F[s]$ is the sum of
$1/q$ over the first $s$ primes $> B$ compatible with a forced set $F \subseteq S$
(quantized conservatively at exact integer scale $S_4 = 10^{18}$: ceiling for
$\gap_L$ and $T_F$, floor for $1/c$ and $\capf(c)$). Note that
$T_F[m_0 + \Delta m] - \sum_{j=1}^{m_0} 1/L_{8{,}136-j}$ is strictly decreasing
in $m_0$, giving an exact upper bound $m_{0,\max}$ and progressively stronger
forced sets on every interval $m_0 \in [m_{0,\mathrm{lo}}, m_{0,\mathrm{hi}}]$.

Moreover, \eqref{eq:delta-budget} immediately yields an elementary per-rank
prime bound: ordering the $d + \Delta m$ primes of $S \setminus L$ in descending
order $x_1 > x_2 > \dots$, for every rank $j \ge \Delta m + 1$, if $x_j > B$
then $x_1,\dots,x_j \in S_{>B}$ each contribute at least $1/B - 1/x_j$ to the
third sum in~\eqref{eq:delta-budget}, forcing
\begin{equation}\label{eq:per-rank}
  x_j < \frac{j B}{j - \Delta m - B\cdot\gap_L}
  \qquad (j \ge \Delta m + 1),
\end{equation}
which also holds trivially when $x_j \le B$. The exact floor $H_j$ strictly
below the right-hand side of~\eqref{eq:per-rank} bounds
$X \le P(n)^{\Delta m}\prod_{j=\Delta m+1}^{\Delta m + d_{\max}} H_j$.

We verify all six $\Delta m$ regimes as summarized in Table~\ref{tab:tier2}:
\begin{itemize}
  \item \textbf{$\Delta m = 3$ ($\omega = 8{,}138$) and $\Delta m = 11$
  ($\omega = 8{,}146$):} Inequality~\eqref{eq:delta-budget} forces
  $|F_3| = 2{,}184$ ($\operatorname{digits}(M_{F_3}) = 12{,}837$) and
  $|F_{11}| = 832$ ($\operatorname{digits}(M_{F_{11}}) = 4{,}362$), while
  \eqref{eq:tail-budget} bounds $d_{\max} = 161$ and $d_{\max} = 291$. One
  witness each for $Q_{\lceil} = \lceil(10^{D-1}-1)/P(L)\rceil$ certifies
  $D = 51{,}798$ ($\Delta m = 3$) and $D = 41{,}895$ ($\Delta m = 11$), while
  one witness each with~\eqref{eq:per-rank} certifies $P(n) > 10^{3{,}685}$ and
  $P(n) > 10^{103}$.

  \item \textbf{$\Delta m = 19$ ($\omega = 8{,}154$):} Here $|F_{19,0}| = 529$
  and $m_{0,\max} = 377$. Splitting $[0, 377]$ into $4$ intervals
  $[0,180]$, $[181,240]$, $[241,300]$, $[301,377]$ raises $|F|$ from $529$ to
  $643, 798, 1{,}201$, refuting $n < 10^{41{,}127}$ ($\ge 41{,}128$ digits) in
  $4$ witnesses.

  \item \textbf{$\Delta m = 27$ ($\omega = 8{,}162$):} Here $|F_{27,0}| = 391$
  ($392$--$1{,}233$ across intervals) and $m_{0,\max} = 448$. Under $n \le 10^{40{,}921} - 1$
  ($Q_{\lceil} = \lceil(10^{40{,}921}-1)/P(L)\rceil$), we partition $[0,448]$
  into $8$ intervals: $[0,120]$, $[121,175]$, $[176,225]$, $[226,270]$,
  $[271,310]$, $[311,345]$, $[346,380]$, $[381,448]$. The first and last two
  intervals close directly ($1$ pair each); the five middle intervals each close
  by a linear forcing cascade over $p \in L \setminus F_{\mathrm{base}}$
  (guarding $p < L_{8{,}136 - m_{0,\mathrm{hi}}}$ so omitted primes remain
  disjoint from the top $m_0$ primes of $L$), totaling $683$ exact witnesses.

  \item \textbf{$\Delta m = 35$ and $\Delta m \ge 43$ (Zero-Float Lagrangian):}
  For any integer $0 < A < B$, define
  $f_A(q) := \log_B(q/B) + A(1/B - 1/q)$. Since $f_A(q) > 0$ for $q > B$ and
  $f_A(q) \le 0$ for $q \le B$, for every nonempty $W_c \subseteq T_c$ we have
  $\sum_{q\in W_c} f_A(q) \ge \sum_{q\in T_c} f_A(q)$. Hence, if every
  non-forced star $c$ satisfies the exact integer-power condition
  \begin{equation}\label{eq:lagrangian}
    c^v B^{v(|T_c|-1) - u_c} \le \Bigl(\prod_{q\in T_c} q\Bigr)^v,
    \qquad u_c := \lfloor v A\,\delta(c)\rfloor,
  \end{equation}
  for some positive integer $v$, then every star term $\sum_{q\in W_c} f_A(q) - f_A(c)$
  is nonnegative, yielding
  \[
    \log_B\!\left(\frac{n}{P(L)}\right) > \Delta m - A\left(\gap_L + \frac{\Delta m}{B}\right)
    = \Delta m\left(1 - \frac{A}{B}\right) - A\cdot\gap_L.
  \]
  With $A_{35} = 6{,}900$ on all $1{,}171$ non-$F_{35}$ stars,
  $r_{35} > 34{,}241/1{,}000$, and the exact integer check
  $P_{\mathrm{top}}^{1000} B^{34{,}241} \ge 10^{1000(40{,}921 - 40{,}729)}$
  (where $P_{\mathrm{top}} = \lfloor P(L)/10^{40{,}729}\rfloor$) proves
  $n > 10^{40{,}921}$ ($\ge 40{,}922$ digits). Globally, $A_{\mathrm{glob}} = 15{,}050 < B$
  passes~\eqref{eq:lagrangian} on all $1{,}455$ stars of $L$; since $1 - A_{\mathrm{glob}}/B > 0$,
  the bound increases with $\Delta m$, and at $\Delta m = 43$ we get
  $r_{43} > 40{,}970/1{,}000$, certifying $n > 10^{40{,}958}$ ($\ge 40{,}959$
  digits) for all $\Delta m \ge 43$.
\end{itemize}
Combining all regimes with Theorems~\ref{thm:extremal} and~\ref{thm:a-ge-2}
completes the proof of Theorem~\ref{thm:tier2}.

\begin{table}[t]
\centering
\caption{Summary of the certified Tier~2 exclusion regimes ($846$ total exact
$(U,R)$ witnesses plus zero-float Lagrangian star checks).}
\label{tab:tier2}
\footnotesize
\setlength{\tabcolsep}{4.5pt}
\begin{tabular}{llcl r}
\toprule
Regime & Method & Forced $|F|$ & Certified bound & Pairs \\
\midrule
$P(n) \le 326{,}056$ & Sound $2$D knapsack + cascade & $104 \to 256$ & $P(n) \ge 326{,}057$ & $155$ \\
$\omega(n) = 8{,}138$ ($\Delta m=3$) & Single-box CRT + rank $H_j$ & $2{,}184$ & $\ge 51{,}798$ digits; $P(n) > 10^{3{,}685}$ & $2$ \\
$\omega(n) = 8{,}146$ ($\Delta m=11$) & Single-box CRT + rank $H_j$ & $832$ & $\ge 41{,}895$ digits; $P(n) > 10^{103}$ & $2$ \\
$\omega(n) = 8{,}154$ ($\Delta m=19$) & $4$ $m_0$-intervals on $[0,377]$ & $529 \to 1{,}201$ & $\ge 41{,}128$ digits & $4$ \\
$\omega(n) = 8{,}162$ ($\Delta m=27$) & $8$ $m_0$-intervals + cascades & $392 \to 1{,}233$ & $\ge 40{,}922$ digits & $683$ \\
$\omega(n) = 8{,}170$ ($\Delta m=35$) & Zero-float Lagrangian ($6{,}900$) & $299$ & $\ge 40{,}922$ digits & $0$ \\
$\omega(n) \ge 8{,}178$ ($\Delta m\ge 43$) & Zero-float Lagrangian ($15{,}050$) & $4$ & $\ge 40{,}959$ digits & $0$ \\
\midrule
Outside Leaf~2 / $a \ge 2$ & Thm.~\ref{thm:extremal} / Thm.~\ref{thm:a-ge-2} & --- & $\ge 44{,}628$ / $\ge 1{,}808{,}611$ digits & $0$ \\
\bottomrule
\end{tabular}
\end{table}

\section{Limitations: The Sharp \texorpdfstring{$r=2$}{r = 2} Giant-Prime Structural Barrier}
\label{sec:barrier}

Theorem~\ref{thm:tier2} shows that any small-$\omega(n)$ counterexample must
possess enormous prime factors, yet CRT + $2$D lattice refutation cannot close
the conjecture by itself. The obstruction is already sharp when only $r = 2$
giant primes complete a base set $S_0$.

\begin{theorem}[$r = 2$ Giant-Prime Structural Barrier]\label{thm:barrier}
Suppose $n = M_0 P_1 P_2$ is an $a = 1$ composite Agoh--Giuga counterexample
with $M_0 = \prod_{p\in S_0} p$ ($S_0 \neq \emptyset$) and two remaining primes
$P_1 < P_2$. Set $N_0 := M_0\bigl(1 - \sum_{p\in S_0} 1/p\bigr) \in \mathbb{Z}$
and $t_0(N_0) := \bigl(1 + \sqrt{1 + 4N_0}\bigr)/(2N_0)$. Then:
\begin{enumerate}
  \item $N_0 \ge 3$ is odd;
  \item $P_1 P_2$ obeys the unconditional strict upper bound
  \begin{equation}\label{eq:r2-bound}
    P_1 P_2 < t_0(N_0)^3 M_0^2 + t_0(N_0)^2 M_0.
  \end{equation}
  In particular, for every $M_0 \ge 10^5$ (hence for every counterexample),
  $P_1 P_2 < 0.4523 M_0^2$.
\end{enumerate}
\end{theorem}

\begin{proof}
Multiplying $1/P_1 + 1/P_2 - 1/(M_0 P_1 P_2) = N_0 / M_0$ by $M_0 P_1 P_2$ yields
$M_0(P_1 + P_2) - 1 = N_0 P_1 P_2$. Since $M_0, P_1, P_2$ are odd, $P_1 + P_2$
is even, so $N_0$ is positive and odd; also $S_0 \neq \emptyset$ gives $N_0 < M_0$.
Because $N_0 P_1 P_2 = M_0 P_1 + M_0 P_2 - 1 > M_0 P_2$, we have
$d_1 := N_0 P_1 - M_0 > 0$, and similarly $d_2 := N_0 P_2 - M_0 > d_1 > 0$,
satisfying
\[
  d_1 d_2 = M_0^2 - N_0
  \qquad\text{and}\qquad
  M_0 P_1 - 1 = (N_0 P_1 - M_0) P_2 = d_1 P_2.
\]
Reducing the second identity modulo $P_2 - 1$, the Carmichael condition
$(P_2 - 1) \mid (M_0 P_1 - 1)$ forces $(P_2 - 1) \mid d_1$, hence
$d_1 \ge P_2 - 1$. If $N_0 = 1$, then $d_1 < d_2 = P_2 - M_0 < P_2 - 1$, a
contradiction; thus $N_0 \ge 3$.

Substituting $P_2 = (M_0 + d_2)/N_0 = (M_0 + (M_0^2 - N_0)/d_1)/N_0$ into
$N_0(d_1 + 1) \ge N_0 P_2$ yields
\[
  f(d_1) := N_0 d_1^2 - (M_0 - N_0) d_1 - (M_0^2 - N_0) \ge 0,
\]
so $d_1 \ge x_+$, where $x_+$ is the unique positive root of $f$. Since
$N_0 \ge 3$ and $M_0 > N_0$, one has $3M_0^2 > 4N_0$, so
$(N_0 - 1)\sqrt{M_0^2 - N_0} \ge 2\sqrt{M_0^2 - N_0} > M_0 > M_0 - N_0$,
whence $f(\sqrt{M_0^2 - N_0}) > 0$ and thus $x_+ < \sqrt{M_0^2 - N_0}$. Because
\[
  P_1 P_2 = \frac{(M_0 + d_1)(M_0 + d_2)}{N_0^2}
          = \frac{2M_0^2 - N_0 + M_0\bigl(d_1 + (M_0^2 - N_0)/d_1\bigr)}{N_0^2}
\]
is strictly decreasing in $d_1$ on $(0, \sqrt{M_0^2 - N_0}]$, we obtain
\[
  P_1 P_2 \le \frac{(M_0 + x_+)(M_0 + (M_0^2 - N_0)/x_+)}{N_0^2}
          = \frac{(M_0 + x_+)(x_+ + 1)}{N_0},
\]
using $N_0(x_+ + 1) = M_0 + (M_0^2 - N_0)/x_+$. Now $t_0 = t_0(N_0)$ satisfies
$N_0 t_0^2 = t_0 + 1$, so $f(t_0 M_0) = N_0(t_0 M_0 + 1) > 0$, implying
$x_+ < t_0 M_0$. Substituting $x_+ < t_0 M_0$ into the strictly increasing
expression above gives
\[
  P_1 P_2 < \frac{(1 + t_0) M_0 (t_0 M_0 + 1)}{N_0}
          = t_0^3 M_0^2 + t_0^2 M_0.
\]
Finally, $t_0(N_0) \le t_0(3) = (1 + \sqrt{13})/6$, where
$t_0(3)^3 = (5 + 2\sqrt{13})/27$ and $t_0(3)^2 < 1$. The exact integer check
$13 \cdot 200{,}000^2 < (27 \cdot 45{,}229 - 500{,}000)^2$ establishes
$t_0(3)^3 + 10^{-5} < 0.4523$, giving $P_1 P_2 < 0.4523 M_0^2$ for all
$M_0 \ge 10^5$.
\end{proof}

\begin{remark}[Why local CRT cannot settle giant-prime completions]
Even if every prime of $L$ were forced, the full CRT modulus
$M_L = P(L)\lambda(L)$ has exact decimal digit length $54{,}909$
($\operatorname{digits}(P(L)) = 40{,}733$, $\operatorname{digits}(\lambda(L)) = 14{,}176$),
whereas $M_0^2 \approx P(L)^2 \approx 10^{81{,}464}$. Because
$t_0(N_0)^3 \sim N_0^{-3/2}$, the inequality $M_L < t_0(N_0)^3 M_0^2$ remains
satisfied across the entire range $3 \le N_0 \lesssim 10^{17{,}704}$. Eliminating
giant-prime completions therefore lies strictly beyond $2$-dimensional lattice
refutation on local CRT congruences. Note also that the $+t_0(N_0)^2 M_0$ term
in~\eqref{eq:r2-bound} is mathematically necessary: at the exact boundary
$d_1 = x_+$ (corresponding to $d_1 = P_2 - 1$), $P_1 P_2$ exceeds
$t_0(N_0)^3 M_0^2$ by $\Theta(M_0)$.
\end{remark}

\section{Reproducibility, Exact Data, and Disclosures}\label{sec:repro}

All computational theorems in this paper are independently reproducible from
scratch using Python~3 standard library alone, with zero external dependencies
and zero floating-point arithmetic in the verifier. Table~\ref{tab:digests}
records the exact canonical integer invariants verified by the standalone
script \path{src/verify_agoh_giuga_certificate.py}.

\begin{table}[htbp]
\centering
\caption{Canonical exact integer invariants certified without floating-point arithmetic.}
\label{tab:digests}
\small
\begin{tabular}{llrr}
\toprule
Quantity & Definition & Digits & Bits \\
\midrule
$P(L)$ & $\prod_{p\in L_{27,692}} p$ ($|L|=8{,}135$, $\max=321{,}053$) & $40{,}733$ & $135{,}311$ \\
$P_{8138}$ & $P(L)\cdot 321{,}073\cdot 321{,}109\cdot 321{,}143$ & $40{,}750$ & $135{,}366$ \\
$M_L$ & $P(L)\cdot \operatorname{lcm}_{p\in L}(p-1)$ & $54{,}909$ & $182{,}403$ \\
\bottomrule
\end{tabular}

\vspace{0.25em}
\footnotesize
Binary big-endian SHA-256 digests:\\
$P(L)$: \texttt{826dad79c2ef410c85a9bd4cb8f3b1ea01dddfabfec8363badc427ca9f4ea7bb}\\
$P_{8138}$: \texttt{54e7192a3b8d0ea5cc571755ff05d9c4e39f788f73acedee41eae134dd439276}\\
$M_L$: \texttt{e37b93406f70df4da505ed071a96a0e9f581ab8c770a0bdefc875e98ed836db0}
\end{table}

\paragraph{Reproduction commands.}
From the project root directory (\texttt{agoh\_giuga/}), the certificate is
generated and independently verified by:
{\small
\begin{verbatim}
python3 src/generate_agoh_giuga_certificate.py
python3 src/verify_agoh_giuga_certificate.py \
  certificates/agoh_giuga_tier2_certificate.json
python3 -m unittest tests/test_agoh_giuga_tier2.py -v
\end{verbatim}
}
On CPython~3.13.5, generation takes approximately $20$ seconds, producing the
$2.68\,\mathrm{MB}$ JSON certificate
(\texttt{agoh\_giuga\_tier2\_certificate.json}), and the standalone zero-float
verifier independently audits all $7$ leaves, the exhaustive mod-$16$ table,
the $a \ge 2$ inequalities, all $846$ lattice witnesses, and the zero-float
Lagrangian checks in approximately $15.6$ seconds. The unit test suite includes
an abstract-syntax-tree audit confirming that the verifier contains zero
floating-point constructs, together with fail-closed tamper tests.

\paragraph{Computational and AI assistance disclosure.}
The mathematical proofs, Star-Pack algorithm, forced-prime CRT cascades, exact
certificate generator, standalone zero-float verifier, and manuscript text were
developed with AI research-assistant tooling. No interactive proof assistant
(such as Lean, Coq, or Isabelle) was used. Every mathematical deduction was
independently audited, and every computational assertion is certified
deterministically by the accompanying zero-float verifier and SHA-256 release
manifest. Author attribution, affiliations, contact details, and acknowledgments
remain anonymized pending user confirmation prior to journal submission.

\bibliographystyle{plain}
\bibliography{references}

\end{document}
